\begin{center}
\begin{tikzpicture}[x=.88cm,y=.88cm,font=\small]
  \fill[paper] (0,0) rectangle (12.3,6.65);

  \draw[source!70,line width=.8pt] (.52,2.02) .. controls (3.5,1.65) and
    (7.15,2.64) .. (11.82,1.72);
  \node[text=source,below,font=\scriptsize] at (8.2,1.85) {长江方向（示意）};
  \draw[source!55,dashed] (.7,4.52) .. controls (4.2,4.12) and
    (7.45,4.92) .. (11.65,4.42);
  \node[text=source,above,font=\scriptsize] at (8.75,4.72) {黄河方向（示意）};

  \fill[dynasty!44] (7.18,4.0) ellipse (1.56 and 1.0);
  \node[align=center,text=white] at (7.18,4.0) {魏\\洛阳与北方仓储};

  \fill[timeline!28] (2.13,2.12) ellipse (1.5 and .92);
  \node[align=center] at (2.13,2.12) {蜀汉\\成都、汉中};

  \fill[source!26] (9.96,1.04) ellipse (1.5 and .76);
  \node[align=center] at (9.96,1.04) {孙吴\\武昌、建业};

  \fill[dynasty!24] (5.0,2.26) ellipse (1.15 and .73);
  \node[align=center,font=\scriptsize] at (5.0,2.26) {荆州、夷陵\\江汉节点};

  \fill[timeline!21] (4.15,5.43) ellipse (1.12 and .66);
  \node[align=center,font=\scriptsize] at (4.15,5.43) {关中、陇右\\山道与粮运};

  \fill[source!21] (10.37,4.95) ellipse (1.14 and .72);
  \node[align=center,font=\scriptsize] at (10.37,4.95) {淮南、合肥\\水陆前线};

  \draw[->,very thick,color=dynasty]
    (5.88,3.77) .. controls (5.32,3.25) and (4.91,2.86) .. (4.63,2.69);
  \node[text=dynasty,below,font=\scriptsize,align=center] at (5.3,3.17)
    {魏、吴争夺\\江汉与江淮};

  \draw[->,very thick,color=timeline]
    (3.35,2.4) .. controls (3.86,2.43) and (4.14,2.4) .. (4.28,2.35);
  \node[text=timeline,below,font=\scriptsize] at (3.78,1.68)
    {221--222年夷陵};

  \draw[->,thick,dashed,color=timeline]
    (2.96,2.83) .. controls (3.2,3.63) and (3.62,4.3) .. (4.01,4.72);
  \node[text=timeline,left,font=\scriptsize,align=right] at (3.18,4.03)
    {汉中—陇右\\北伐方向};

  \draw[->,thick,dashed,color=source]
    (8.5,3.77) .. controls (9.08,4.2) and (9.36,4.55) .. (9.48,4.69);
  \node[text=source,right,font=\scriptsize,align=left] at (8.78,4.47)
    {魏、吴的\\淮南争夺};

  \draw[->,thick,color=dynasty]
    (8.14,3.23) .. controls (8.84,2.6) and (9.22,1.97) .. (9.43,1.7);
  \node[text=dynasty,right,font=\scriptsize,align=left] at (8.33,2.48)
    {长江水军\\与行政腹地};

  \node[draw=source!75,fill=paper,rounded corners=2pt,align=center,
    inner sep=3pt,font=\scriptsize] at (7.02,.62)
    {三国的都城与名号并未消除道路、水面和山口造成的限制；\\
    这些节点反复决定动员能否抵达前线};
\end{tikzpicture}

\smallskip
\small\textit{图\quad 魏、蜀汉、孙吴的政治中心与主要竞争走廊，约220--229年：据《三国志·文帝纪》《先主传》《吴主传》及魏、蜀、吴相关列传，并参照《资治通鉴》卷六十九至七十一，概括洛阳、成都、建业（武昌）、荆州和汉中之间的关系。政权名称对应不同时点的中心安排；色块不等于疆界，箭线只提示史料所见的战略方向，不表示完整战线、行军路线或可通行程度。}
\end{center}
